\chapter{Netzwerk-Spiel}
\label{chapter:netzwerkspiel}
% \begin{figure}[H]
%     \centering
%     \includegraphics[width=.7\textwidth]{Grundlagen_Info/07_Netzwerke/Skript/Appendices/Zettel.pdf}
%     \caption{Zettel für das Netzwerk-Spiel}
% \end{figure}
% \clearpage

In diesem Spiel schicken Sie sich in Gruppen gegenseitig Nachrichten zu und übernehmen dabei selbst die Aufgaben der Schichten. Jede Gruppe besteht aus einem Router und drei oder vier Computern. Am Anfang hat jeder Computer nur eine feste Geräte-Nummer, die \gls{MAC}-Adresse. Eine \gls{IP}-Adresse (Adresse im Netz) bekommt er erst vom Router. Die Druckvorlagen (Anleitungen für Computer und Router, Router-Tabellen, Gerätekarten, Ethernet-Frames, IP-Pakete und Nachrichtenzettel) liegen im Ordner \texttt{Netzwerkspiel} dieses Themas.

\begin{table}[H]
	\centering
	\small
	\begin{tabular}{p{4cm}p{4.2cm}p{5.6cm}}
		Spielmaterial & Schicht & Inhalt \\\midrule
		Nachrichtenzettel & Anwendungsschicht & Nachricht, höchstens 10 Zeichen \\
		(Kopf des Zettels) & Transportschicht & Sequenznummer (\enquote{Seq 2 von 3}) \\
		IP-Kärtchen & Internetschicht & Von-\acrshort{IP}, An-\acrshort{IP}, TTL \\
		Ethernet-Frame (A4) & Netzzugangsschicht & Ziel-\acrshort{MAC} und Quell-\acrshort{MAC} \\
	\end{tabular}
	\caption{Spielmaterial und Schichten im Netzwerk-Spiel}
\end{table}

\begin{enumerate}
	\item \textbf{Runde 1 (IP-Adressen verteilen und zustellen):} Die Computer holen beim Router eine \acrshort{IP}-Adresse. Der Router schreibt \acrshort{MAC}- und \acrshort{IP}-Adresse in seine Tabelle. Jeder Computer teilt eine Frage auf mehrere Nachrichtenzettel auf, nummeriert sie, steckt jeden Zettel in ein IP-Paket und gibt es dem Router. Der Router sucht die \acrshort{MAC}-Adresse zur Ziel-\acrshort{IP}, steckt das IP-Paket in das passende Ethernet-Frame und legt es in die Tischmitte. Er stellt die Pakete absichtlich nicht immer in der Reihenfolge zu, in der sie ankommen.
	\item \textbf{Runde 2 (der Router vergisst):} Der Router weiss nicht mehr, welche \acrshort{MAC}-Adresse zu welcher \acrshort{IP}-Adresse gehört. Er fragt laut \enquote{Wer hat die IP-Adresse \dots?} und schreibt die Antworten in eine neue Tabelle.
	\item \textbf{Runde 3 (zwei Gruppen verbunden):} Pakete an die andere Gruppe gibt der Router in einem neuen Frame an den anderen Router weiter, der sie wieder in einem neuen Frame zustellt. Auf dem IP-Paket steht eine Zahl TTL (\enquote{Time to live}). Jeder Router zieht 1 ab, bei 0 wird das Paket weggeworfen.
\end{enumerate}

Der Router ist streng: Stimmt bei einem Paket etwas nicht (z.\,B. ungültige Adresse, Empfänger unbekannt, TTL abgelaufen), wirft er es weg und sagt dem Absender, weshalb. Im echten Internet wird ein Paket oft auch ohne Meldung weggeworfen.

\begin{myremark}
	Das Spiel vereinfacht einiges. Im echten Netz verpackt der Computer selbst das IP-Paket in ein Frame. Die \acrshort{MAC}-Adresse des Empfängers erfragt er wie der Router in Runde 2, und innerhalb eines Netzes leitet ein Switch das Frame weiter. Ein Router packt ein Paket erst dann neu ein, wenn es in ein anderes Netz geht, wie in Runde 3. \acrshort{MAC}-Adressen haben ausserdem 6 statt 3 Bytes.
\end{myremark}

Für die Nachbesprechung eignen sich diese Fragen: Welche Angabe gehört zu welcher Schicht? Was hat sich in Runde 3 auf dem Weg geändert, was nicht? Weshalb braucht es zwei Adressen (\acrshort{MAC} und \acrshort{IP})? Wo entstand ein Stau, und wozu dient die Sequenznummer?

